\begin{abstract}
O desenvolvimento de software com IA deixou de ser uma questão de prompts isolados e passou a se organizar em frameworks que estruturam o processo com artefatos persistentes, papéis agênticos, execução integrada e validação. Faltava, porém, uma revisão que tomasse esses frameworks operacionais como objeto, com uma lente de Engenharia de Software, em vez de tratar arquiteturas de agentes de forma abstrata ou técnicas isoladas do ciclo de vida. Este artigo apresenta uma revisão sistemática de literatura conduzida sob o protocolo de Kitchenham, que reuniu 37 estudos publicados entre 2022 e 2026 após triagem em duas passadas e três rodadas de snowballing, incluindo apenas artigos acadêmicos e preprints arquivados com registro verificável. A síntese temática em modo híbrido derivou seis temas indutivos e os confrontou com uma taxonomia proposta de seis dimensões (especificação, contexto, papéis, execução, validação e portabilidade). Os resultados sustentam empiricamente o deslocamento do prompt para o processo, presente em 34 dos 37 estudos, e validam cinco das seis dimensões; portabilidade é fraca e quase sempre diagnóstica, e dois eixos emergem fora da taxonomia, a colaboração humano-IA e a segurança. Identificamos três eixos de tensão recorrentes e mostramos que contexto e rastreabilidade, e não a geração de código em si, são os gargalos práticos transversais. Concluímos com uma agenda de pesquisa orientada a processo, centrada em benchmarks que avaliem o pipeline, comparação direta entre frameworks e avaliação longitudinal de governança.
\end{abstract}

\noindent\textbf{Palavras-chave:} desenvolvimento de software com IA; agentes de IA; revisão sistemática; engenharia de software; frameworks agênticos.
